\documentclass[11pt,a4paper]{article}

\usepackage[margin=1in]{geometry}
\usepackage{graphicx}
\usepackage{listings}
\usepackage{xcolor}
\usepackage{booktabs}
\usepackage{longtable}
\usepackage{hyperref}
\usepackage{parskip}
\usepackage{fancyhdr}
\usepackage{titlesec}
\usepackage{amsmath}

\definecolor{codebg}{RGB}{245,245,245}
\definecolor{codekw}{RGB}{0,0,180}
\definecolor{codecomment}{RGB}{100,100,100}
\definecolor{codestr}{RGB}{163,21,21}

\lstset{
  backgroundcolor=\color{codebg},
  basicstyle=\ttfamily\small,
  keywordstyle=\color{codekw}\bfseries,
  commentstyle=\color{codecomment}\itshape,
  stringstyle=\color{codestr},
  breaklines=true,
  frame=single,
  language=Python,
  numbers=left,
  numberstyle=\tiny,
  xleftmargin=2em,
  showstringspaces=false
}

\pagestyle{fancy}
\fancyhf{}
\rhead{ICS1512 --- ML Algorithms Lab}
\lhead{Experiment 1}
\cfoot{\thepage}

\title{\textbf{Experiment 1: Working with Python Packages --- \\ NumPy, SciPy, Scikit-Learn, Matplotlib}}
\author{
  M.Tech (Integrated) Computer Science \& Engineering \\
  ICS1512 --- Machine Learning Algorithms Laboratory \\
  Sri Sivasubramaniya Nadar College of Engineering, Chennai
}
\date{Academic Year 2026--2027 (Odd) \quad Batch: 2024--2029}

\begin{document}
\maketitle

\section*{Aim}
To explore the core functions of the Python libraries \textbf{NumPy}, \textbf{Pandas}, \textbf{SciPy},
\textbf{Scikit-learn}, and \textbf{Matplotlib}; to download benchmark datasets from public repositories
(UCI/Kaggle) and identify the appropriate machine learning model for each; and to carry out the complete
machine learning workflow --- data loading, exploratory data analysis, preprocessing, feature selection,
data splitting, and performance evaluation --- on five representative datasets.

\section{Library Exploration}

\subsection{NumPy}
NumPy provides the \texttt{ndarray} object for fast, vectorized numerical computation. Key operations
explored include array creation (\texttt{zeros}, \texttt{ones}, \texttt{arange}, \texttt{linspace}),
indexing/slicing, reshaping, element-wise arithmetic with broadcasting, matrix operations
(\texttt{dot}, transpose, determinant, inverse), boolean masking, and random number generation.

\begin{lstlisting}
a = np.array([1, 2, 3, 4, 5])
mat = np.arange(12).reshape(3, 4)
print("Sum:", a.sum(), "Mean:", a.mean(), "Std:", a.std())
print(np.dot(m1, m2))          # matrix multiplication
print(np.linalg.inv(m1))       # matrix inverse
\end{lstlisting}

\subsection{Pandas}
Pandas builds on NumPy to provide labeled \texttt{Series} and \texttt{DataFrame} structures for
tabular data. Operations explored include DataFrame creation, inspection (\texttt{info}, \texttt{describe}),
boolean filtering, handling missing values with \texttt{fillna}, and aggregation via \texttt{groupby}.

\begin{lstlisting}
df = pd.DataFrame({"Name": [...], "Age": [...], "Marks": [...]})
df2.fillna(df2["Marks"].mean())
df.groupby("Result")["Marks"].mean()
\end{lstlisting}

\subsection{SciPy}
SciPy extends NumPy with modules for statistics (\texttt{scipy.stats}), special functions
(\texttt{scipy.special}), and linear algebra (\texttt{scipy.linalg}). Descriptive statistics
(mean, median, mode, skewness, kurtosis), a one-sample $t$-test, the Gamma function, and solving a
linear system $Ax = b$ were demonstrated.

\begin{lstlisting}
t_stat, p_val = stats.ttest_1samp(data, popmean=15)
x = linalg.solve(A, b)
\end{lstlisting}

\subsection{Matplotlib}
Matplotlib was used to generate line plots, scatter plots, bar charts, and histograms in a single
$2\times2$ subplot grid, illustrating the core visualization primitives used throughout EDA
(Figure~\ref{fig:matplotlib}).

\begin{figure}[h]
  \centering
  \includegraphics[width=0.85\linewidth]{figs/fig_01.png}
  \caption{Matplotlib: line, scatter, bar, and histogram plots.}
  \label{fig:matplotlib}
\end{figure}

\subsection{Scikit-learn}
Scikit-learn provides a uniform \texttt{fit}/\texttt{predict}/\texttt{transform} API for the ML
workflow. A baseline tour --- loading the Iris dataset, splitting it, scaling with
\texttt{StandardScaler}, and fitting a \texttt{LogisticRegression} classifier --- was carried out to
illustrate this pattern before applying it across all five datasets.

\section{Dataset Exploration and ML Task Identification}

Public datasets analogous to those in the UCI Machine Learning Repository and Kaggle were used. Since
Iris, the Diabetes-progression dataset, and 8$\times$8 handwritten digit images (a lightweight stand-in
for MNIST) are bundled directly with Scikit-learn, they were loaded via
\texttt{sklearn.datasets}. The Loan Amount Prediction and Email Spam datasets are not bundled with
Scikit-learn, so structurally faithful synthetic datasets (same feature schema as the Kaggle/UCI
originals) were generated in code; in an environment with open internet access these loading cells
should be replaced with \texttt{pd.read\_csv(...)} on the downloaded CSV files.

\subsection{Iris Dataset --- Classification}
Four numeric features (sepal/petal length and width) are used to classify flowers into three species.
Figure~\ref{fig:iris} shows the sepal-length distribution by species and the petal length--width
scatter plot, which shows near-perfect linear separability between classes.

\begin{figure}[h]
  \centering
  \includegraphics[width=0.9\linewidth]{figs/fig_02.png}
  \caption{Iris EDA: sepal length distribution and petal length vs.\ width by species.}
  \label{fig:iris}
\end{figure}

An ANOVA F-test (\texttt{SelectKBest}) ranked petal length and petal width as the most discriminative
features. A \texttt{DecisionTreeClassifier} trained on an 80/20 split achieved high accuracy, confirmed
by the confusion matrix in Figure~\ref{fig:iris_cm}.

\begin{figure}[h]
  \centering
  \includegraphics[width=0.5\linewidth]{figs/fig_04.png}
  \caption{Iris --- Decision Tree confusion matrix on the held-out test set.}
  \label{fig:iris_cm}
\end{figure}

\subsection{Loan Amount Prediction --- Regression}
The target variable (loan amount, a continuous quantity) makes this a \textbf{regression} problem.
Missing values in \texttt{LoanAmount} and \texttt{CreditHistory} were imputed (median/most-frequent),
categorical fields (\texttt{Education}, \texttt{PropertyArea}) were one-hot encoded, and numeric
fields were standardized inside a \texttt{ColumnTransformer} pipeline. Figure~\ref{fig:loan} shows the
loan-amount distribution and its spread across property-area categories.

\begin{figure}[h]
  \centering
  \includegraphics[width=0.9\linewidth]{figs/fig_05.png}
  \caption{Loan dataset EDA: loan amount distribution and boxplot by property area.}
  \label{fig:loan}
\end{figure}

A \texttt{LinearRegression} model was fit on the preprocessed features and evaluated using MAE and
$R^2$, both reported directly in the notebook output.

\subsection{Predicting Diabetes --- Classification}
The Scikit-learn diabetes-progression target (a continuous disease-progression score) was binarized
around its median to reproduce the classic diabetes-risk \textbf{classification} task. A Chi-square
test (after min--max scaling to non-negative values) ranked feature importance, and a
\texttt{RandomForestClassifier} was trained and evaluated. Figure~\ref{fig:diab} shows the feature
correlation heatmap and BMI distribution by class.

\begin{figure}[h]
  \centering
  \includegraphics[width=0.9\linewidth]{figs/fig_06.png}
  \includegraphics[width=0.45\linewidth]{figs/fig_07.png}
  \caption{Diabetes dataset: feature correlation heatmap and BMI by diabetic class.}
  \label{fig:diab}
\end{figure}

\subsection{Email Spam Classification --- Classification}
Word-frequency-style numeric features (frequency of ``free'', ``win'', ``offer'', capital-run length,
number of links) were used to classify messages as spam or not spam, mirroring the UCI Spambase
schema. Figure~\ref{fig:spam} shows the clear separation in the frequency of the word ``free'' and in
capital-run length between spam and non-spam classes. An ANOVA F-test ranked feature relevance, and a
\texttt{GaussianNB} classifier was trained on standardized features.

\begin{figure}[h]
  \centering
  \includegraphics[width=0.9\linewidth]{figs/fig_08.png}
  \caption{Spam dataset EDA: word-frequency and capital-run-length distributions by class.}
  \label{fig:spam}
\end{figure}

\subsection{Handwritten Character Recognition (MNIST-style) --- Classification}
The Scikit-learn \texttt{load\_digits} dataset (8$\times$8 grayscale images of digits 0--9,
1{,}797 samples) was used as a lightweight stand-in for the full 28$\times$28 MNIST dataset.
Figure~\ref{fig:digits} shows sample digit images. Pixel intensities were standardized and a
\texttt{Support Vector Machine} (RBF kernel) was trained, achieving high classification accuracy; the
confusion matrix is shown in Figure~\ref{fig:digits_cm}.

\begin{figure}[h]
  \centering
  \includegraphics[width=0.9\linewidth]{figs/fig_09.png}
  \caption{Sample handwritten digit images from the dataset.}
  \label{fig:digits}
\end{figure}

\begin{figure}[h]
  \centering
  \includegraphics[width=0.55\linewidth]{figs/fig_10.png}
  \caption{Handwritten digits --- SVM confusion matrix on the held-out test set.}
  \label{fig:digits_cm}
\end{figure}

\section{Summary Table}

\begin{longtable}{p{3.2cm} p{2.8cm} p{3.3cm} p{4cm}}
\toprule
\textbf{Dataset} & \textbf{Type of ML Task} & \textbf{Feature Selection Technique} & \textbf{Suitable ML Algorithm} \\
\midrule
\endhead
Iris Dataset & Classification (Supervised) & ANOVA F-test (SelectKBest) & Decision Tree / Logistic Regression \\
Loan Amount Prediction & Regression (Supervised) & Correlation / domain-based selection & Linear Regression \\
Predicting Diabetes & Classification (Supervised) & Chi-square test & Random Forest Classifier \\
Classification of Email Spam & Classification (Supervised) & ANOVA F-test (SelectKBest) & Naive Bayes / Logistic Regression \\
Handwritten Character Recognition / MNIST & Classification (Supervised) & Variance / PCA on pixel intensities & Support Vector Machine (SVM) / CNN \\
\bottomrule
\end{longtable}

\section{Inference and Learning Outcomes}
\begin{itemize}
  \item \textbf{NumPy} underlies all numerical computation in the ML stack through fast, vectorized
        array operations, broadcasting, and linear algebra routines.
  \item \textbf{Pandas} provides labeled tabular structures with built-in tools for cleaning, filtering,
        grouping, and merging real-world data.
  \item \textbf{SciPy} extends NumPy with statistical tests and scientific routines needed for
        hypothesis testing and numerical solving.
  \item \textbf{Matplotlib/Seaborn} make exploratory visualization possible, revealing distribution
        shape, class separability, outliers, and feature correlations before modeling.
  \item \textbf{Scikit-learn} unifies preprocessing, feature selection, model fitting, and evaluation
        behind a single consistent API.
  \item Across the five datasets, tasks spanned \textbf{multi-class classification} (Iris, Digits,
        Spam, Diabetes) and \textbf{regression} (Loan Amount), demonstrating how task type governs the
        choice of algorithm, evaluation metric (accuracy/F1 vs.\ MAE/$R^2$), and feature-selection
        method (ANOVA/Chi-square for classification, correlation-based methods for regression).
  \item Careful \textbf{preprocessing} (imputation, encoding, scaling) and a proper
        \textbf{train/test split} were essential in every case for obtaining reliable, generalizable
        performance estimates.
\end{itemize}

\end{document}
